% !TeX root = ../thuthesis-example.tex

\begin{acknowledgements}
  三年求学时光荏苒，转眼即将迎来下一关键阶段。值此硕士毕业之际，我想向所有在我求学路上给予我帮助和支持的师长、同学、家人们表达我最诚挚的感谢。

  衷心感谢我的导师罗禹贡老师一直以来的全方位指导。罗老师治学严谨务实，学术功底深厚，始终以身作则。罗老师始终密切关注我的学业情况与科研进展，每一次的学术交流后鞭辟入里的问题分析与每一次论文修改时细致入微的方向点拨，都让我受益匪浅。罗老师的指点和教诲将激励我在未来的道路上砥砺前行。

  衷心感谢高博麟老师在我硕士期间的言传身教与鼎力支持。高老师对科研问题的敏锐洞见、产业落地的魄力与对汽车的热情始终是我科研的引航灯。除科研外，高老师始终关心并持续给予我许多试错与提升的重要机会。后续博士求学期间我一定继续努力，沉淀出经得起产业考验的硬技术。衷心感谢钟薇老师在我硕士期间的悉心指导与无微不至的帮助，钟老师总是耐心地解答我的疑惑，在我面对压力退缩时给予我宝贵的建议和鼓励。同时，感谢李克强老师给我继续攻读博士学位的机会；感谢李升波老师、温福喜老师、何雷老师、许庆老师、曲小波老师、王红老师、王凯老师、徐少兵老师等在硕士期间给予我各方面的指导和帮助。

  感谢课题组全体师长和同学给予我的热心帮助和支持。感谢刘勃老师、张放老师、吕尧老师、肖雅月老师、卢彦博老师、王彤老师、王勇老师、王越老师、蔡孟池老师的指导和支持；感谢韩硕、刘家熙、刘鸿祥、王璐瑶、杨皓宇、万科科、林立、程浩、石佳、关书睿、李鹏飞、王宙、熊宣威、张傲、黎铸新、吕子垣等师兄师姐在科研中的帮助；感谢徐熙标、王怡然两位在Onsite的帮助；感谢何瑞坤、邹恒多、黄梓谦、邱逸凡、扶尚宇、王艺航、王京坤、关婷文等同学在天门山比赛中的大力支持，我们是最棒的团队；感谢龚林豪、王坤鹏两位师弟在东风项目不遗余力的付出；感谢肖壮锐、葛亚敏、吕敏、闵刚峰等同学在项目上的支持；感谢黄建程、柴祺等好友交流关照，以及在清华有幸结识的钱康安、郑梓昂、武淑敏、任晗啸、赵享、王以诺、王立坤、赵亦悦等一众好友，唯愿友谊长存。

  最后，要特别感谢我的父亲、母亲以及所有家家人们，你们时刻教导我要乐观积极，不畏艰难险阻，给予我家庭的温暖和最无私的爱；感谢我的女友张佳，遇见你是我人生中最幸运的事，感谢你和崽子家族最真挚的爱与陪伴支持，让我一次又一次勇敢向前，未来我们携崽子家族全体继续攻坚克难勇攀高峰！

  本课题得到国家重点研发计划“机场飞行区有人无人作业装备混合运行智能管控技术”(2024YFB2605200) 的支持，特此致谢。

\end{acknowledgements}
